% !PRINTED 271
\weberpage{271}
Zur Erreichung dieses Zieles gen\"ugt es auch hier, wenn $n$ keinen quadratischen Teiler hat, da\ss{} $\Phi_{a,0,\partial}$, und wenn $n$ eine Primzahl ist, wenn $\Phi_{1,0,n}$ f\"ur $q=0$ endlich bleibt.

Bei diesen Rechnungen machen wir Gebrauch von den Entwickelungen:
\weq{6}{\begin{aligned}
uv &= q^{-\frac{n+1}{24}}\Pi(1+q^{n(2h+1)})(1+q^{2h+1}),\\
u_1v_1 &= q^{-\frac{n+1}{24}}\Pi(1-q^{n(2h+1)})(1-q^{2h+1}),\\
u_2v_2 &= 2q^{\frac{n+1}{12}}\Pi(1+q^{2hn})(1+q^{2h}),
\end{aligned}}
woraus durch Entwickelung nach Potenzen $q$:
\weq{7}{\begin{aligned}
uv &= q^{-\frac{n+1}{24}}(1+q+q^3+q^4+q^5+q^6+q^7+2q^8\dots),\\
u_1v_1 &= q^{-\frac{n+1}{24}}(1-q-q^3+q^4-q^5+q^6-q^7+2q^8\dots),\\
u_2v_2 &= 2q^{\frac{n+1}{12}}(1+q^2+q^4+2q^6+2q^8+\cdots).
\end{aligned}}
Die letzteren Formeln sind richtig f\"ur $n>7$ (f\"ur $n=3,5,7$ sind die Glieder von $q^3, q^5, q^7$ an zu modifizieren). F\"ur $n=7$, $n=23$ zeigen diese Ausdr\"ucke, da\ss{} $A$ selbst f\"ur $q=0$ endlich bleibt, woraus f\"ur diese F\"alle die Gleichungen folgen:
\weq{8}{\begin{aligned}
n&=7, & A&=0,\\
n&=23, & A&=1.
\end{aligned}}
Durch einfache Rechnung findet man ferner noch:
\weq{9}{\begin{aligned}
n&=31, & (A^2-B)^2-A&=0,\\
n&=47, & A^2-A-B&=2,\\
n&=71, & A^3-4A^2+2A-B&=1.
\end{aligned}}
Nicht ganz so einfach gestaltet sich die Rechnung f\"ur ein zusammengesetztes $n$. So mu\ss{} man z.\,B. f\"ur $n=15$ die Bedingung der Endlichkeit f\"ur $q=0$ nicht nur f\"ur $\Phi_{1,0,15}$, sondern auch f\"ur $\Phi_{3,0,5}$ ber\"ucksichtigen. Diese beiden aber gen\"ugen. Man erh\"alt so:
\weq{10}{n=15, \quad A^3-AB+1=0.}
2. Ist $n \equiv 3$ (mod 8), so sind dieselben Schl\"usse zu ziehen, wenn wir setzen:
\weq{11}{\begin{aligned}
4A &= u^2v^2 - u_1^2v_1^2 - u_2^2v_2^2,\\
B &= u^2v^2u_1^2v_1^2 + u^2v^2u_2^2v_2^2 - u_1^2v_1^2u_2^2v_2^2,
\end{aligned}}
%CONT>>
